\documentclass[11pt]{article}

% ---------- Packages ----------
\usepackage[margin=1in]{geometry}
\usepackage{amsmath, amssymb}   % math
\usepackage{graphicx}           % \includegraphics
\usepackage{booktabs}           % nicer tables
\usepackage{float}              % [H] placement
\usepackage[font=small, skip=6pt]{caption}
\usepackage{subcaption}         % side-by-side scope captures
\usepackage{xcolor}
\usepackage{listings}        % code blocks, pure TeX (no shell-escape)
\usepackage[hidelinks]{hyperref}

% Draft marker: red text for anything still to be filled in from the lab notebook.
\newcommand{\fillin}[1]{\textcolor{red}{[#1]}}

\graphicspath{{pictures/}}

% ---------- Code block style (applies to every code block) ----------
% AVR assembly: listings ships no dialect for it, so define one.
\lstdefinelanguage{avrasm}{
  morekeywords={
    add,adc,adiw,sub,subi,sbc,sbci,sbiw,and,andi,or,ori,eor,com,neg,inc,dec,tst,clr,ser,
    mov,movw,ldi,ld,ldd,lds,st,std,sts,push,pop,in,out,lpm,
    cp,cpc,cpi,cpse,breq,brne,brcs,brcc,brsh,brlo,brmi,brpl,brge,brlt,brts,brtc,
    rjmp,jmp,ijmp,rcall,call,icall,ret,reti,nop,sleep,wdr,
    lsl,lsr,rol,ror,asr,swap,sbi,cbi,sbic,sbis,sbrc,sbrs,sbr,cbr,set,clt,
    sei,cli,clz,sez,clc,sec,
  },
  morekeywords=[2]{.cseg,.dseg,.eseg,.org,.byte,.db,.dw,.equ,.def,.include,.macro,.endmacro},
  alsoletter={.},          % so directives like .cseg tokenize as one word
  morecomment=[l]{;},
  morestring=[b]",
  sensitive=false,
}
\lstdefinestyle{labcode}{
  basicstyle=\ttfamily\small,
  numbers=left, numberstyle=\tiny\color{gray!70}, numbersep=9pt,
  xleftmargin=2.2em, framexleftmargin=2.2em,   % background covers the number gutter
  breaklines=true, breakatwhitespace=true,
  frame=lines, framesep=2mm,
  backgroundcolor=\color{gray!8},
  keywordstyle=\color{blue!65!black}\bfseries,
  keywordstyle=[2]\color{purple!70!black},
  commentstyle=\color{green!40!black},
  stringstyle=\color{orange!85!black},
  showstringspaces=false,
  tabsize=4,
  captionpos=t,
}
\lstset{style=labcode}
\renewcommand{\lstlistingname}{Code}   % caption says "Code 1" instead of "Listing 1"

% ---------- Report info (edit these) ----------
\newcommand{\course}{ECE 3218L}
\newcommand{\labtitle}{Lab 6: The Timer Unit}
\newcommand{\studentname}{Daniel Xu}
\newcommand{\labdate}{10/15/2026}

% article.cls always typesets title -> author -> date, so the date goes in the
% title block to sit directly under it. \date{} is required; empty means \today.
\title{\labtitle \\[0.3em] \large \labdate}
\author{\studentname \\ \small \course{} \,}
\date{}

% One oscilloscope capture inside a two-column figure.
\newcommand{\scope}[3]{%
  \begin{subfigure}[t]{0.48\textwidth}
    \centering
    \includegraphics[width=\linewidth]{#1}
    \caption{#2}
    \label{#3}
  \end{subfigure}}

\begin{document}
\maketitle

\section{Overview}
% Timer1: 16-bit counter, output compare (OC1A on PB1) and input capture (ICP1 on PB0).
% What the three programs do: rectangular wave, a sequence of tones, half-frequency copy of an input.

\section{Procedure}
% Connections: OC1A (PB1, Uno pin 9) to the scope / piezo, ICP1 (PB0, Uno pin 8) from the
% function generator, common ground. Generator checked for 0-5 V before connecting.

% \begin{figure}[H]
%   \centering
%   \includegraphics[width=0.45\textwidth]{connections.png}
%   \caption{Connections.}
%   \label{fig:connections}
% \end{figure}

\begin{enumerate}
\item \textbf{Program 1: Rectangular wave.}
% SETWAVE: inputs (pw, duty code, prescaler), high/low widths from right shifts,
% 75% and 87.5% as the inverse of 25% and 12.5%, CTC + toggle OC1A, ISR swaps OCR1A.
% Prescaler choice (pre-lab questions a-i): count = fcc * pw / sp, sp = 1, count = 36364.

\begin{table}[H]
  \centering
  \caption{Program 1 at 440 Hz ($pw = 36364$ clicks, prescaler 1).}
  \label{tab:duty}
  \begin{tabular}{crrrr}
    \toprule
    & \multicolumn{2}{c}{Expected} & \multicolumn{2}{c}{Measured} \\
    \cmidrule(lr){2-3} \cmidrule(l){4-5}
    Duty cycle & High & Low & High & Low \\
    \midrule
    12.5\% & 284.1~$\mu$s & 1.989~ms & \fillin{} & \fillin{} \\
    25\%   & 568.2~$\mu$s & 1.705~ms & \fillin{} & \fillin{} \\
    50\%   & 1.136~ms     & 1.136~ms & \fillin{} & \fillin{} \\
    75\%   & 1.705~ms     & 568.2~$\mu$s & \fillin{} & \fillin{} \\
    87.5\% & 1.989~ms     & 284.1~$\mu$s & \fillin{} & \fillin{} \\
    \bottomrule
  \end{tabular}
\end{table}

\item \textbf{Program 2: Tones.}
% Tone periods copied to SRAM, then for each: SETWAVE at 50%, wait 500 ms, stop the clock
% (clear TCCR1B), wait 500 ms. What the scale sounded like.

\begin{table}[H]
  \centering
  \caption{Tone periods at prescaler 1 ($pw = 16~\text{MHz} / f$).}
  \label{tab:tones}
  \begin{tabular}{lrr}
    \toprule
    Note & Frequency (Hz) & $pw$ (clicks) \\
    \midrule
    C4 & 261.63 & 61155 \\
    D4 & 293.66 & 54485 \\
    E4 & 329.63 & 48539 \\
    F4 & 349.23 & 45815 \\
    G4 & 392.00 & 40816 \\
    A4 & 440.00 & 36364 \\
    B4 & 493.88 & 32397 \\
    C5 & 523.25 & 30578 \\
    D5 & 587.33 & 27242 \\
    E5 & 659.25 & 24270 \\
    \bottomrule
  \end{tabular}
\end{table}

\item \textbf{Program 3: Input capture.}
% Timer1 free-running at clk/8 (0.5 us per click). Capture ISR: PERIOD = ICR1 - previous.
% Compare ISR: OCR1A += PERIOD, toggles OC1A once per input period -> half the input frequency.

\begin{table}[H]
  \centering
  \caption{Program 3 input and output frequencies.}
  \label{tab:capture}
  \begin{tabular}{rrr}
    \toprule
    Input (Hz) & Expected output (Hz) & Measured output (Hz) \\
    \midrule
    1000 & 500 & \fillin{} \\
    \fillin{} & \fillin{} & \fillin{} \\
    \fillin{} & \fillin{} & \fillin{} \\
    \bottomrule
  \end{tabular}
\end{table}
\end{enumerate}

\section{Discussion}
\subsection{Program 1: Accuracy and Limits}
% Did the scope match Table 1? Any deviation, and why (or why not).
% Limits of the program and of the high/low portions: 16-bit OCR1A, ISR must finish
% before the next match (~40 cycles), pw >= 8 for 12.5%. Note on 90% vs 87.5%.

\subsection{Program 2: Frequency Range}
% Smallest and largest frequency: 16 MHz / (sp * 65535) at the low end,
% ISR time at the high end. Where each limit comes from.

\subsection{Program 3: Measuring a Period}
% Results explained to the TA. Rollover handling (16-bit subtraction).
% Longest measurable period: 65535 * 0.5 us = 32.77 ms at clk/8.

\section{Conclusion}
% Two or three sentences.

\appendix
\section{Full Code}
\label{app:code}

% Read straight from the Lab6 project so the listings always match the code.
% \footnotesize so long comment lines fit.
\lstinputlisting[language=avrasm,basicstyle=\ttfamily\footnotesize,caption={Program 1. Rectangular wave (main.asm).},label={app:main}]{../../Lab6/main.asm}

\lstinputlisting[language=avrasm,basicstyle=\ttfamily\footnotesize,caption={SETWAVE subroutine and compare ISR, used by Programs 1 and 2 (wave.inc).},label={app:wave}]{../../Lab6/wave.inc}

\lstinputlisting[language=avrasm,basicstyle=\ttfamily\footnotesize,caption={Program 2. Tones (tones.asm).},label={app:tones}]{../../Lab6/tones.asm}

\lstinputlisting[language=avrasm,basicstyle=\ttfamily\footnotesize,caption={Program 3. Input capture (capture.asm).},label={app:capture}]{../../Lab6/capture.asm}

\end{document}
